\section{Discussion}
\subsection{Why Does SVD Work for Attention?}
The success of SVD-Lin hinges on the fact that the attention matrix $A = \text{softmax}(QK^\top/\sqrt{d})$ is often approximately low-rank. This property has been observed empirically: the singular values of the attention matrix decay rapidly, especially in lower layers. By compressing $K$ and $V$ using the left singular vectors of $H$, we are effectively projecting the entire sequence into a subspace that captures most of the variance. This is analogous to principal component analysis (PCA) applied to the token representations.

\subsection{Trade-Offs and Practical Recommendations}
\begin{itemize}
    \item For tasks with short sequences ($N \le 128$), rank 128 provides almost no accuracy loss and a 1.6× speedup – highly recommended.
    \item For resource-constrained environments (mobile, edge), rank 64 offers excellent efficiency with acceptable accuracy drop.
    \item For long sequences ($N > 512$), rank 128 is still sufficient; the linear scaling becomes a major advantage.
\end{itemize}

\subsection{Limitations}
\begin{enumerate}
    \item \textbf{SVD overhead}: Although randomized SVD is fast, it still adds some cost. For very small batches, this overhead can dominate. However, in typical training scenarios with batch size $\ge 16$, the benefit outweighs the cost.
    \item \textbf{Not end-to-end differentiable in the rank}: The discrete choice of rank is a hyperparameter; we cannot learn it via gradient descent easily.
    \item \textbf{MPS support}: As of PyTorch 2.0, `svd\_lowrank` is not implemented on MPS, forcing a CPU→GPU copy. This will be resolved in future versions.
\end{enumerate}

\subsection{Future Work}
\subsubsection{Adaptive Rank Selection}
Instead of a fixed rank, we can choose $r$ dynamically based on the singular value decay. For example, set $r$ such that $\sigma_{r+1} / \sigma_1 < \epsilon$. This would automatically allocate more capacity to sequences that need it.

\subsubsection{Per-Head SVD}
Different attention heads may capture different linguistic properties. Computing separate $U_r$ per head could improve performance, though at higher cost.

\subsubsection{Incremental SVD for Causal Attention}
For autoregressive generation, we can update the SVD incrementally as new tokens arrive, using rank-1 updates. This would make SVD-Lin suitable for streaming settings.

\subsubsection{Combination with Pruning}
SVD compression can be combined with weight pruning to further reduce model size. For instance, after training, we could apply SVD to the feed-forward network weights as well.